% !TEX program = tectonic
% =============================================================================
%  AULA 5 — Redes sem Fio Emergentes  (atualizada a 2026)
%  Curso: Redes sem Fio  |  Profa. Anelise Munaretto
%  Wi-Fi 6/6E/7, LoRa/Sigfox, 5G NR, satélite LEO (Starlink), DTN, FANET, VANET.
%  Compilar: tectonic aula5.tex
% =============================================================================
\documentclass[aspectratio=169,11pt]{beamer}
% =============================================================================
%  PREÁMBULO COMÚN para as aulas de Redes sem Fio (Beamer + TikZ)
%  Incluido com % =============================================================================
%  PREÁMBULO COMÚN para as aulas de Redes sem Fio (Beamer + TikZ)
%  Incluido com % =============================================================================
%  PREÁMBULO COMÚN para as aulas de Redes sem Fio (Beamer + TikZ)
%  Incluido com \input{preamble.tex} en cada aula.
%  Paleta de cores e estilos são definidos aqui uma única vez.
% =============================================================================

\usepackage[T1]{fontenc}
\usepackage{amsmath,amssymb,mathtools}
\usepackage{graphicx}
\usepackage{tikz}
\usetikzlibrary{positioning,arrows.meta,calc,backgrounds,decorations.pathmorphing,fit,shadows,shapes.symbols,shapes.geometric}
\usepackage{pgfplots}
\pgfplotsset{compat=1.16}
\usepackage{tabularx,booktabs,multirow,array}
\usepackage[normalem]{ulem}
\usepackage{hyperref}

% ---------- Paleta de cores própria ----------
\definecolor{aneliseAzul}{RGB}{20,60,110}
\definecolor{aneliseVerde}{RGB}{30,140,70}
\definecolor{aneliseLaranja}{RGB}{215,120,20}
\definecolor{aneliseGris}{RGB}{90,90,90}

% ---------- Tema ----------
\usetheme{Madrid}
\usecolortheme{whale}
\setbeamercolor{structure}{fg=aneliseAzul}
\setbeamercolor{title}{fg=aneliseAzul}
\setbeamercolor{frametitle}{fg=aneliseAzul}
\setbeamercolor{block title}{fg=aneliseAzul,bg=aneliseGris!12!white}
\setbeamercolor{block body}{bg=aneliseGris!7!white}
\hypersetup{colorlinks=true,linkcolor=aneliseAzul,urlcolor=aneliseVerde}

% ---------- Comandos auxiliares ----------
\newcommand{\kurose}{Kurose \& Ross, \emph{Computer Networking: a Top-Down Approach}}
\newcommand{\forouzan}{Forouzan, \emph{Data Communications and Networking}}
\newcommand{\stallings}{Stallings, \emph{Wireless Communications \& Networks}}
\newcommand{\tanenbaum}{Tanenbaum, \emph{Computer Networks}}
\newcommand{\rfc}[1]{\href{https://www.rfc-editor.org/rfc/rfc#1.txt}{RFC~#1}}
% -----------------------------------------------------------------------------
%  Estilos TikZ comuns para desenhar redes (posicionamento relativo)
% -----------------------------------------------------------------------------
\tikzset{
  host/.style={draw,rounded corners=2mm,fill=aneliseAzul!15,inner sep=1.5mm,font=\scriptsize},
  rt/.style={draw,rounded corners=1mm,fill=aneliseLaranja!25,inner sep=1mm,font=\scriptsize},
  nube/.style={draw,cloud,aspect=2.2,fill=aneliseGris!15,inner sep=1.5mm,font=\scriptsize},
  el/.style={draw,rounded corners=1.2mm,fill=aneliseGris!18,inner sep=0.8mm,font=\scriptsize},
  enl/.style={thick},
  enlA/.style={thick,-Stealth},
} en cada aula.
%  Paleta de cores e estilos são definidos aqui uma única vez.
% =============================================================================

\usepackage[T1]{fontenc}
\usepackage{amsmath,amssymb,mathtools}
\usepackage{graphicx}
\usepackage{tikz}
\usetikzlibrary{positioning,arrows.meta,calc,backgrounds,decorations.pathmorphing,fit,shadows,shapes.symbols,shapes.geometric}
\usepackage{pgfplots}
\pgfplotsset{compat=1.16}
\usepackage{tabularx,booktabs,multirow,array}
\usepackage[normalem]{ulem}
\usepackage{hyperref}

% ---------- Paleta de cores própria ----------
\definecolor{aneliseAzul}{RGB}{20,60,110}
\definecolor{aneliseVerde}{RGB}{30,140,70}
\definecolor{aneliseLaranja}{RGB}{215,120,20}
\definecolor{aneliseGris}{RGB}{90,90,90}

% ---------- Tema ----------
\usetheme{Madrid}
\usecolortheme{whale}
\setbeamercolor{structure}{fg=aneliseAzul}
\setbeamercolor{title}{fg=aneliseAzul}
\setbeamercolor{frametitle}{fg=aneliseAzul}
\setbeamercolor{block title}{fg=aneliseAzul,bg=aneliseGris!12!white}
\setbeamercolor{block body}{bg=aneliseGris!7!white}
\hypersetup{colorlinks=true,linkcolor=aneliseAzul,urlcolor=aneliseVerde}

% ---------- Comandos auxiliares ----------
\newcommand{\kurose}{Kurose \& Ross, \emph{Computer Networking: a Top-Down Approach}}
\newcommand{\forouzan}{Forouzan, \emph{Data Communications and Networking}}
\newcommand{\stallings}{Stallings, \emph{Wireless Communications \& Networks}}
\newcommand{\tanenbaum}{Tanenbaum, \emph{Computer Networks}}
\newcommand{\rfc}[1]{\href{https://www.rfc-editor.org/rfc/rfc#1.txt}{RFC~#1}}
% -----------------------------------------------------------------------------
%  Estilos TikZ comuns para desenhar redes (posicionamento relativo)
% -----------------------------------------------------------------------------
\tikzset{
  host/.style={draw,rounded corners=2mm,fill=aneliseAzul!15,inner sep=1.5mm,font=\scriptsize},
  rt/.style={draw,rounded corners=1mm,fill=aneliseLaranja!25,inner sep=1mm,font=\scriptsize},
  nube/.style={draw,cloud,aspect=2.2,fill=aneliseGris!15,inner sep=1.5mm,font=\scriptsize},
  el/.style={draw,rounded corners=1.2mm,fill=aneliseGris!18,inner sep=0.8mm,font=\scriptsize},
  enl/.style={thick},
  enlA/.style={thick,-Stealth},
} en cada aula.
%  Paleta de cores e estilos são definidos aqui uma única vez.
% =============================================================================

\usepackage[T1]{fontenc}
\usepackage{amsmath,amssymb,mathtools}
\usepackage{graphicx}
\usepackage{tikz}
\usetikzlibrary{positioning,arrows.meta,calc,backgrounds,decorations.pathmorphing,fit,shadows,shapes.symbols,shapes.geometric}
\usepackage{pgfplots}
\pgfplotsset{compat=1.16}
\usepackage{tabularx,booktabs,multirow,array}
\usepackage[normalem]{ulem}
\usepackage{hyperref}

% ---------- Paleta de cores própria ----------
\definecolor{aneliseAzul}{RGB}{20,60,110}
\definecolor{aneliseVerde}{RGB}{30,140,70}
\definecolor{aneliseLaranja}{RGB}{215,120,20}
\definecolor{aneliseGris}{RGB}{90,90,90}

% ---------- Tema ----------
\usetheme{Madrid}
\usecolortheme{whale}
\setbeamercolor{structure}{fg=aneliseAzul}
\setbeamercolor{title}{fg=aneliseAzul}
\setbeamercolor{frametitle}{fg=aneliseAzul}
\setbeamercolor{block title}{fg=aneliseAzul,bg=aneliseGris!12!white}
\setbeamercolor{block body}{bg=aneliseGris!7!white}
\hypersetup{colorlinks=true,linkcolor=aneliseAzul,urlcolor=aneliseVerde}

% ---------- Comandos auxiliares ----------
\newcommand{\kurose}{Kurose \& Ross, \emph{Computer Networking: a Top-Down Approach}}
\newcommand{\forouzan}{Forouzan, \emph{Data Communications and Networking}}
\newcommand{\stallings}{Stallings, \emph{Wireless Communications \& Networks}}
\newcommand{\tanenbaum}{Tanenbaum, \emph{Computer Networks}}
\newcommand{\rfc}[1]{\href{https://www.rfc-editor.org/rfc/rfc#1.txt}{RFC~#1}}
% -----------------------------------------------------------------------------
%  Estilos TikZ comuns para desenhar redes (posicionamento relativo)
% -----------------------------------------------------------------------------
\tikzset{
  host/.style={draw,rounded corners=2mm,fill=aneliseAzul!15,inner sep=1.5mm,font=\scriptsize},
  rt/.style={draw,rounded corners=1mm,fill=aneliseLaranja!25,inner sep=1mm,font=\scriptsize},
  nube/.style={draw,cloud,aspect=2.2,fill=aneliseGris!15,inner sep=1.5mm,font=\scriptsize},
  el/.style={draw,rounded corners=1.2mm,fill=aneliseGris!18,inner sep=0.8mm,font=\scriptsize},
  enl/.style={thick},
  enlA/.style={thick,-Stealth},
}

\title[Redes sem Fio]{Redes sem Fio Emergentes}
\subtitle{Aula 5}
\author[Anelise Munaretto]{Profa. Anelise Munaretto}
\institute[Curso]{Redes sem Fio --- Ciência da Computação}
\date{2026}

\begin{document}

\begin{frame}[plain]
  \titlepage
  \begin{center}
    \scriptsize Conteúdo atualizado 2026: 802.11ax/be, LPWAN (LoRa/Sigfox), 5G, NTN.
  \end{center}
\end{frame}

\begin{frame}{Objetivos de aprendizagem}
  Ao final desta aula, você será capaz de:
  \begin{itemize}
    \item \textbf{Comparar} as gerações do IEEE 802.11 (ax/be) e \textbf{explicar} o papel do \emph{multi-link operation} (MLO) e dos canais de 320~MHz;
    \item \textbf{Calcular} o tempo no ar de pacotes LoRaWAN a partir de SF, BW e CR;
    \item \textbf{Relacionar} os casos de uso do 5G (eMBB, uRLLC, mMTC) às tecnologias-chave;
    \item \textbf{Descrever} a arquitetura de constelações LEO e o papel dos enlaces a laser ISL;
    \item \textbf{Analisar} arquiteturas emergentes (DTN, FANET, VANET) quanto a atraso, mobilidade e topologia.
  \end{itemize}
  \begin{alertblock}{Pré-requisitos}
    Aulas 2 (camada física: modulação, SNR, sensibilidade) e 4 (acesso ao meio: MAC, CSMA/CA, RTS/CTS, 802.11e).
  \end{alertblock}
\end{frame}

\section{IEEE 802.11 (Wi-Fi)}

\begin{frame}{LAN sem fio IEEE 802.11: evolução}
  \begin{center}
    \scriptsize
    \begin{tabular}{@{}llll@{}}
      \toprule
      \textbf{Padrão} & \textbf{Banda} & \textbf{Taxa} & \textbf{Notas}\\
      \midrule
      802.11b & 2,4~GHz & 11~Mbps & DSSS, chip de 11\\
      802.11a & 5--6~GHz & 54~Mbps & OFDM (52 subport.)\\
      802.11g & 2,4--5~GHz & 54~Mbps & OFDM\\
      802.11n (Wi-Fi 4) & 2,4/5~GHz & 600~Mbps & MIMO 4$\times$4, OFDM 108\\
      802.11ac (Wi-Fi 5) & 5~GHz & 6,9~Gbps & MIMO 8$\times$8, QAM-256\\
      802.11ax (Wi-Fi 6) & 2,4/5/6~GHz & 9,6~Gbps & OFDMA, QAM-1024, MU-MIMO\\
      802.11be (Wi-Fi 7) & 2,4/5/6~GHz & $\sim$46~Gbps & 320~MHz, QAM-4096, mMLS\\
      \bottomrule
    \end{tabular}
  \end{center}
  \begin{itemize}
    \item todos usam \textbf{CSMA/CA} para acesso múltiplo
    \item todos têm versão de estação base e ad hoc
  \end{itemize}
\end{frame}

\begin{frame}{802.11: canais e associação}
  \begin{columns}
    \begin{column}{0.55\textwidth}
      \textbf{802.11b}: espectro 2,4--2,485~GHz dividido em 11 canais
      \begin{itemize}
        \item o admin do AP escolhe a frequência do AP
        \item interferência possível: o canal pode ser o mesmo que o do AP vizinho!
      \end{itemize}
      \textbf{Associação do hospedeiro}
      \begin{enumerate}
        \item varre os canais, escutando \textbf{beacons} (SSID e MAC)
        \item seleciona um AP para se associar
        \item pode fazer autenticação (WPA3)
        \item normalmente executa DHCP para obter IP na sub-rede do AP
      \end{enumerate}
    \end{column}
    \begin{column}{0.45\textwidth}
      \begin{center}
        \begin{tikzpicture}[font=\scriptsize,
            ap/.style={draw,circle,minimum size=7mm,fill=aneliseLaranja!30,inner sep=0.5mm},
            n/.style={draw,circle,minimum size=6mm,fill=aneliseAzul!15,inner sep=0.4mm},
          ]
          \node[ap](ap1){AP1};\node[ap,right=2.2cm of ap1](ap2){AP2};
          \node[n,below=1.2cm of ap1](a1){a};\node[n,below=1.2cm of ap2](a2){b};
          \node[n,below=2.2cm of ap1](a3){c};\node[n,below=2.2cm of ap2](a4){d};
          \draw[thick,aneliseLaranja] (ap1) -- (a1);
          \draw[thick,aneliseLaranja] (ap2) -- (a2);
          \draw[dashed,thick,aneliseGris] (ap1) -- (a3);
          \draw[dashed,thick,aneliseGris] (ap2) -- (a4);
          \draw[dashed,thick,aneliseLaranja!60] (ap1) to[bend left=20] node[midway,above=2mm,font=\tiny]{canal 6} (ap2);
        \end{tikzpicture}\\[1mm]
        {\scriptsize \textbf{Figura 1:} Canais não sobrepostos (1, 6 e 11) e associação de estações aos APs.}
      \end{center}
      {\scriptsize Objetivo: canais 1, 6 e 11 não sobrepostos (2,4~GHz).}
    \end{column}
  \end{columns}
\end{frame}

\begin{frame}{Wi-Fi 6 (802.11ax) e Wi-Fi 6E}
  \begin{columns}
    \begin{column}{0.5\textwidth}
      \textbf{Wi-Fi 6 (802.11ax, 2019)}
      \begin{itemize}
        \item \textbf{OFDMA}: subportadoras alocadas a múltiplos usuários no mesmo slot
        \item \textbf{MU-MIMO} uplink/downlink
        \item QAM-1024
        \item objetivo: eficiência em ambientes densos (estádios, aeroportos)
      \end{itemize}
    \end{column}
    \begin{column}{0.5\textwidth}
      \textbf{Wi-Fi 6E (2021)}
      \begin{itemize}
        \item adiciona a \textbf{banda de 6~GHz} (5925--7125~MHz)
        \item mais largura de banda disponível, menos interferência
        \item requer \textbf{WPA3} obrigatório
      \end{itemize}
    \end{column}
  \end{columns}
  \vspace{0.2cm}
  \begin{block}{Wi-Fi 7 (802.11be, 2024)}
    \begin{itemize}
      \item canais de \textbf{320~MHz}, QAM-4096, \emph{multi-link operation} (MLO)
      \item throughput teórico até \textbf{46~Gbps} (em ambientes ideais)
      \item dirigido a streaming 8K, AR/VR, TSN/indústria
    \end{itemize}
  \end{block}
  \note{Enfatizar que o MLO do Wi-Fi 7 usa simultaneamente as bandas de 2,4/5/6 GHz, reduzindo latência e aumentando confiabilidade; o Wi-Fi 6E foi o primeiro a habilitar a banda de 6 GHz.}
\end{frame}

\section{WPAN: Bluetooth e outras}

\begin{frame}{802.15.1: rede de área pessoal (Bluetooth)}
  \begin{columns}
    \begin{column}{0.5\textwidth}
      \begin{itemize}
        \item menos de 10~m de diâmetro
        \item substituta de cabos (mouse, teclado, fones)
        \item ad hoc: sem infraestrutura
        \item mestre/escravos:
          \begin{itemize}
            \item escravos pedem permissão para enviar (ao mestre)
            \item o mestre concede as solicitações
          \end{itemize}
      \end{itemize}
    \end{column}
    \begin{column}{0.5\textwidth}
      \begin{center}
        \begin{tikzpicture}[font=\scriptsize,
            m/.style={draw,circle,minimum size=7mm,fill=aneliseLaranja!35,inner sep=0.5mm},
            s/.style={draw,circle,minimum size=6mm,fill=aneliseAzul!15,inner sep=0.4mm},
            p/.style={draw,circle,minimum size=6mm,fill=aneliseGris!20,inner sep=0.4mm},
          ]
          \node[m](M){M};
          \node[s,above=1.2cm of M](s1){S};
          \node[s,left=1.2cm of M](s2){S};
          \node[s,right=1.2cm of M](s3){S};
          \node[p,below=1.2cm of M](p1){P};
          \draw[thick,aneliseLaranja] (M) -- (s1);
          \draw[thick,aneliseLaranja] (M) -- (s2);
          \draw[thick,aneliseLaranja] (M) -- (s3);
          \draw[dashed,thick,aneliseGris] (M) -- (p1);
          \node[draw,fill=aneliseGris!12,rounded corners=2mm,below=5mm of p1,text width=3.4cm,align=center](leg){M mestre, S escravo,\\P estacionado (inativo)};
        \end{tikzpicture}\\[1mm]
        {\scriptsize \textbf{Figura 2:} Piconet Bluetooth (mestre, escravos, parado).}
      \end{center}
      {\scriptsize Banda 2,4--2,5~GHz; até 721~kbps (BT 1.x); a evolução 802.15.1 até 3~Mbps.}
    \end{column}
  \end{columns}
  \begin{exampleblock}{Evolução (2026)}
    Bluetooth 5.4 (LE Audio, Mesh), UWB (802.15.4z, precisão em cm), Zigbee 3.0, Thread/Matter.
  \end{exampleblock}
\end{frame}

\section{LPWAN: LoRa e Sigfox}

\begin{frame}{LPWAN para IoT: LoRa vs.\ Sigfox}
  \begin{columns}
    \begin{column}{0.5\textwidth}
      \textbf{LoRa / LoRaWAN}
      \begin{itemize}
        \item modulação \textbf{Chirp Spread Spectrum (CSS)}
        \item banda ISM sub-GHz (868~MHz EU, 915~MHz US/BR)
        \item taxas 0,3--50~kbps; alcance 2--15~km
        \item topologia \emph{star-of-stars}: sensores $\rightarrow$ gateways $\rightarrow$ servidor
        \item ALOHA puro + canais com \emph{duty cycle}
        \item autenticação AES-128 (join E2E)
      \end{itemize}
    \end{column}
    \begin{column}{0.5\textwidth}
      \textbf{Sigfox}
      \begin{itemize}
        \item banda ultra-estreita (UNB, 100~Hz)
        \item uplink 100 bps; 140 mensagens/dia
        \item alcance até 50~km (rural)
        \item rede proprietária (operador regional)
      \end{itemize}
      \begin{block}{Comparação prática}
        LoRaWAN: mais flexível, redes privadas, bidirecional.
        Sigfox: ultra-simples, muito baixo consumo, unidirecional.
      \end{block}
    \end{column}
  \end{columns}
  \note{Relacionar SF e tempo no ar: SF maior = símbolos mais longos, maior sensibilidade e alcance, porém menor taxa e maior consumo — base para o exercício resolvido.}
\end{frame}

\begin{frame}{NB-IoT e eMTC (celular LPWA)}
  \begin{itemize}
    \item padrões 3GPP para IoT massivo dentro de LTE/5G:
      \begin{itemize}
        \item \textbf{NB-IoT} (Narrowband IoT): 200~kHz, até 250~kbps, muito boa penetração indoor
        \item \textbf{LTE-M / eMTC}: 1~MHz, até 1~Mbps, suporta voz e mobilidade
      \end{itemize}
    \item comparados com LoRa: maior confiabilidade, QoS garantida, porém maior consumo e custo de módulo
  \end{itemize}
  \begin{center}
    \begin{tikzpicture}[font=\scriptsize,
        c/.style={draw,rounded corners=2mm,fill=aneliseAzul!13,inner sep=1mm,text width=3.2cm,align=center},
      ]
      \node[c](a){LoRaWAN\\(css sub-GHz, sem licença)};
      \node[c,right=2.8cm of a](b){NB-IoT\\(LTE, licença)};
      \node[c,right=2.8cm of b](c1){5G mMTC\\(NR, licença)};
      \draw[<->,thick,aneliseGris] (a.east) -- (b.west);
      \draw[<->,thick,aneliseGris] (b.east) -- (c1.west);
      \node[below=2mm of b,font=\tiny]{alcance, consumo, custo};
    \end{tikzpicture}\\[1mm]
    {\scriptsize \textbf{Figura 3:} Comparação qualitativa LoRaWAN, NB-IoT e 5G mMTC quanto a alcance, consumo e custo.}
  \end{center}
\end{frame}

\section{5G e redes móveis}

\begin{frame}{5G New Radio (NR)}
  \begin{columns}
    \begin{column}{0.5\textwidth}
      \textbf{Três casos de uso}
      \begin{itemize}
        \item \textbf{eMBB} (Enhanced Mobile Broadband): $>$ 1~Gbps, 8K/VR
        \item \textbf{uRLLC} (Ultra-Reliable Low Latency): 1~ms, indústria, condução autônoma
        \item \textbf{mMTC} (Massive Machine-Type): 1M dispositivos/km², IoT massivo
      \end{itemize}
    \end{column}
    \begin{column}{0.5\textwidth}
      \textbf{Tecnologias-chave}
      \begin{itemize}
        \item \emph{beamforming} massivo (Massive MIMO)
        \item \emph{network slicing} (redes virtuais por serviço)
        \item \emph{edge computing} (MEC)
        \item ondas milimétricas (mmWave 26--28~GHz, 60~GHz)
      \end{itemize}
    \end{column}
  \end{columns}
  \begin{exampleblock}{Evolução}
    5G-Advanced (3GPP Rel.\ 18, 2024): IA/ML na rede, RedCap. O 6G já está em pesquisa (3GPP Rel.\ 21, por volta de 2030).
  \end{exampleblock}
\end{frame}

\section{Redes ad hoc, DTN, drones e veículos}

\begin{frame}{Arquiteturas emergentes}
  \begin{columns}
    \begin{column}{0.5\textwidth}
      \textbf{Topologias}
      \begin{itemize}
        \item Ad hoc, infraestrutura, Mesh
        \item Sensores (WSN)
        \item VANET (veículos), FANET (drones)
        \item DTN (Delay/Disruption Tolerant Networks)
        \item Green Networking, SDR, SDN
      \end{itemize}
    \end{column}
    \begin{column}{0.5\textwidth}
      \textbf{DTN (Delay/Disruption Tolerant Networks)}
      \begin{itemize}
        \item tolerância a atrasos e interrupções (espaço, móvel rural)
        \item \emph{store-and-forward} com o protocolo \textbf{Bundle Protocol (RFC 5050)}
      \end{itemize}
      \textbf{Flying Networks (drones)}
      \begin{itemize}
        \item FANET: enxames de UAV, roteamento 3D, mobilidade rápida
      \end{itemize}
      \textbf{Redes veiculares}
      \begin{itemize}
        \item C-V2X (celular) vs.\ 802.11p/ITS-G5
        \item comunicação V2V, V2I, V2P
      \end{itemize}
    \end{column}
  \end{columns}
\end{frame}

\section{Satélite e NTN}

\begin{frame}{Satélite: redes não terrestres (NTN)}
  \begin{columns}
    \begin{column}{0.5\textwidth}
      \textbf{Satélite tradicional}
      \begin{itemize}
        \item GEO (35\,786~km): cobertura ampla, alta latência (125~ms)
        \item MEO/LEO: menor latência
      \end{itemize}
      \textbf{LEO massivo (2026)}
      \begin{itemize}
        \item \textbf{Starlink} (SpaceX): $>$ 7000 satélites, constelação em LEO $\sim$550~km
        \item \textbf{OneWeb, Kuiper, Telesat}
        \item latência 20--40~ms; throughput 100--300~Mbps por terminal
        \item laser ISL (inter-satellite links); integração 5G NTN (3GPP Rel.\ 17+)
      \end{itemize}
    \end{column}
    \begin{column}{0.5\textwidth}
      \begin{center}
        \begin{tikzpicture}[font=\scriptsize]
          \node[draw,fill=aneliseGris!20,rounded corners=2mm](t){terminal terrestre};
          \node[draw,circle,fill=aneliseLaranja!30,above=1.6cm of t](s1){sat LEO};
          \node[draw,circle,fill=aneliseLaranja!30,above=1.3cm of s1](s2){sat LEO};
          \node[draw,circle,fill=aneliseLaranja!30,above=1.3cm of s2](s3){sat LEO};
          \node[draw,fill=aneliseGris!20,above=1.6cm of s3](g){gateway};
          \draw[thick,aneliseVerde] (t) -- (s1);
          \draw[thick,aneliseLaranja] (s1) -- (s2) -- (s3);
          \draw[thick,aneliseVerde] (s3) -- (g);
          \node[right=1mm of s2,font=\tiny]{enlaces a laser ISL};
        \end{tikzpicture}\\[1mm]
        {\scriptsize \textbf{Figura 4:} Constelação LEO com enlaces a laser ISL entre satélites e gateway terrestre.}
      \end{center}
    \end{column}
  \end{columns}
  \note{Contrastar latência: GEO ~125 ms vs.\ LEO 20--40 ms; destacar que a retransmissão entre satélites via ISL a laser permite cobertura global com poucos gateways.}
\end{frame}

\section{SDN, rádio cognitivo e pesquisa}

\begin{frame}{Rádio cognitivo e SDN}
  \begin{columns}
    \begin{column}{0.5\textwidth}
      \textbf{Rádio cognitivo (IEEE 802.22 / TVWS)}
      \begin{itemize}
        \item acesso oportunista ao espectro (escanear, detectar, usar)
        \item \emph{cognitive radio networks}: espectro dinâmico
      \end{itemize}
      \textbf{SDR (Software Defined Radio)}
      \begin{itemize}
        \item processamento do sinal por software; agilidade espectral
      \end{itemize}
    \end{column}
    \begin{column}{0.5\textwidth}
      \textbf{SDN (Software Defined Networking)}
      \begin{itemize}
        \item plano de controle centralizado
        \item aplicado a WMN, VANET, IoT
        \item vantagem: programabilidade da rede
      \end{itemize}
    \end{column}
  \end{columns}
\end{frame}

\begin{frame}{Projetos de pesquisa em redes sem fio}
  \textbf{Áreas principais}: otimização de recursos e protocolos de roteamento.
  \begin{enumerate}
    \item Redes complexas $\times$ dados dinâmicos $\times$ Big Data $\times$ IA
    \item IoT e Smart Cities: LoRaONE com CSMA
    \item Redes ad hoc de drones $\times$ gestão de mobilidade $\times$ otimização de recursos
    \item Inteligência coletiva $\times$ roteamento $\times$ redes sem fio oportunistas $\times$ cenários urbanos
    \item Redes veiculares $\times$ mobilidade urbana $\times$ roteamento
    \item Otimização de recursos $\times$ rádios cognitivos
    \item Internet usando a rede de transporte público de Curitiba
  \end{enumerate}
  \begin{exampleblock}{Renovação}
    Incorporar: \emph{mesh Wi-Fi 6/7 doméstico}, \emph{network slicing 5G}, \emph{TPU/GPU sensing}, \emph{dados abertos de Curitiba} (URBS).
  \end{exampleblock}
\end{frame}

\begin{frame}{Exercício resolvido: LoRaWAN --- tempo no ar}
  \small
  \textbf{Dados}: $SF=12$, $BW=125$~kHz, $CR=4/5$, payload $PL=20$ bytes, header explícito ($H=0$), CRC ativado, $DE=0$.
  \begin{block}{Passo 1 --- Símbolo e preâmbulo}
    $T_{sym}=2^{SF}/BW=4096/125\,000=32{,}768$~ms; \quad $T_{preamble}=12{,}25\cdot T_{sym}\approx 401{,}4$~ms
  \end{block}
  \begin{block}{Passo 2 --- Símbolos de payload}
    $n_{payload}=8+\big\lceil(8PL-4SF+28+16\,CRC-20H)/(4(SF-2DE))\big\rceil\cdot(CR+4)$\\
    $\phantom{n_{payload}}=8+\lceil 156/48\rceil\cdot 8=8+4\cdot 8=40$ símbolos
  \end{block}
  \begin{block}{Passo 3 --- Tempo no ar}
    $T_{payload}=40\cdot 32{,}768\approx 1310{,}7$~ms \quad $\Rightarrow$ \quad $T_{total}=T_{preamble}+T_{payload}\approx \mathbf{1{,}71}$~\textbf{s}
  \end{block}
  \begin{alertblock}{Interpretação}
    Com \emph{duty cycle} de 1\% (EU868), entre duas transmissões deve haver $\geq 100\times 1{,}71\approx 171$~s ($\approx 2{,}9$~min).
  \end{alertblock}
\end{frame}

\begin{frame}{Exercício resolvido: LoRaWAN --- alcance por orçamento de enlace}
  \textbf{Dados}: $P_{TX}=14$~dBm, sensibilidade $S_{RX}=-137$~dBm (SF12/BW125), frequência $868$~MHz.
  \begin{block}{Orçamento de enlace}
    \begin{itemize}
      \item perda máxima admissível: $L_{max}=P_{TX}-S_{RX}=14-(-137)=\mathbf{151}$~\textbf{dB}
      \item perda no espaço livre (modelo): $PL(d)=32{,}44+20\log_{10}(f[\text{MHz}])+20\log_{10}(d[\text{km}])$
      \item a $10$~km: $PL=32{,}44+20\log_{10}(868)+20\log_{10}(10)=32{,}44+58{,}77+20\approx 111$~dB
    \end{itemize}
  \end{block}
  \begin{exampleblock}{Resultado}
    Margem de $151-111\approx 40$~dB a $10$~km: o link é viável mesmo com perdas por obstáculos/multipercurso; na prática, o alcance LoRaWAN fica entre 2--15~km (rural até $\sim$50~km no Sigfox).
  \end{exampleblock}
\end{frame}

\begin{frame}{Resumo da aula}
  \begin{itemize}
    \item \textbf{Wi-Fi 6/7}: OFDMA, MU-MIMO, QAM-1024/4096, canais de 320~MHz e \emph{multi-link operation} (MLO);
    \item \textbf{LPWAN}: LoRaWAN (CSS, \emph{star-of-stars}, \emph{duty cycle}) e Sigfox (UNB) para IoT de baixo custo e consumo; NB-IoT/LTE-M como alternativa celular com QoS;
    \item \textbf{5G NR}: casos eMBB, uRLLC e mMTC viabilizados por \emph{beamforming} massivo, \emph{network slicing} e MEC;
    \item \textbf{Satélites LEO (Starlink)}: latência 20--40~ms, enlaces a laser ISL e integração NTN (3GPP Rel.\ 17+);
    \item \textbf{Arquiteturas emergentes}: DTN (Bundle Protocol), FANET/VANET, SDR e SDN como objetos de pesquisa ativa.
  \end{itemize}
\end{frame}

\begin{frame}{Autoavaliação}
  \small
  \begin{enumerate}
    \item Qual geração do Wi-Fi introduz o \emph{multi-link operation} (MLO)?
      \begin{itemize}
        \item (a) 802.11n (Wi-Fi 4)
        \item (b) 802.11ac (Wi-Fi 5)
        \item (c) 802.11be (Wi-Fi 7)
      \end{itemize}
      \pause
      \textbf{Gabarito:} (c)
    \item Em LoRaWAN, aumentar o spreading factor (SF) de 7 para 12:
      \begin{itemize}
        \item (a) aumenta a taxa de dados e reduz o alcance
        \item (b) aumenta o alcance e o tempo no ar, reduzindo a taxa
        \item (c) não altera o alcance nem a taxa
      \end{itemize}
      \pause
      \textbf{Gabarito:} (b)
    \item Qual caso de uso do 5G exige latência da ordem de 1~ms?
      \begin{itemize}
        \item (a) mMTC
        \item (b) eMBB
        \item (c) uRLLC
      \end{itemize}
      \pause
      \textbf{Gabarito:} (c)
  \end{enumerate}
\end{frame}

\begin{frame}{Laboratório sugerido: LoRaWAN na prática ou simulação LPWAN}
  \begin{block}{Opção 1 --- Node LoRaWAN com The Things Network (TTN)}
    \begin{itemize}
      \item hardware: microcontrolador (ESP32/Arduino) + módulo LoRa (ex.: SX1276/RFM95), banda 868/915~MHz;
      \item registrar o dispositivo na TTN (OTAA), enviar payloads e observar RSSI/SNR no console;
      \item variar SF (7$\rightarrow$12) e medir tempo no ar e consumo; comparar com o exercício resolvido.
    \end{itemize}
  \end{block}
  \begin{block}{Opção 2 --- Simulação LPWAN no ns-3}
    \begin{itemize}
      \item usar o módulo LoRaWAN do ns-3 (ou o LoraWanSim) com gateways e end-devices;
      \item variar SF, tamanho de payload e \emph{duty cycle}; coletar taxa de entrega, atraso e energia;
      \item avaliar o impacto do \emph{spreading factor} na cobertura e no consumo.
    \end{itemize}
  \end{block}
\end{frame}

\section{Referências}

\begin{frame}{Referências}
  \begin{small}
  \begin{itemize}
    \item Green Networking: \url{http://dx.doi.org/10.1109/SURV.2011.113010.00106}
    \item Radio Cognitive: \url{http://dx.doi.org/10.1109/98.788210}
    \item SDN: \url{http://dx.doi.org/10.1109/MCOM.2013.6553676}
    \item Flying Ad Hoc Network: \url{http://dx.doi.org/10.1109/TVT.2015.2414819}
    \item Dados Abertos Curitiba: \url{http://www.curitiba.pr.gov.br/dadosabertos/}
    \item URBS Itibus: \url{https://www.urbs.curitiba.pr.gov.br/mobile/itibus}
    \item Complex networks: \url{https://arxiv.org/pdf/cond-mat/0303516.pdf}
    \item IEEE 802.11be (Wi-Fi 7): \url{https://arxiv.org/pdf/2007.13401.pdf}
  \end{itemize}
  \end{small}
  \begin{block}{Novas referências (2026)}
    \begin{itemize}
      \item 3GPP Rel.\ 17/18 NTN e RedCap
      \item IEEE 802.11ax/be (IEEE Std 802.11ax-2021, 802.11be-2024)
      \item LoRaWAN 1.0.4 (LoRa Alliance, 2020)
    \end{itemize}
  \end{block}
\end{frame}

\begin{frame}{Leitura recomendada}
  \begin{itemize}
    \item \textbf{\kurose{}} --- Capítulo 7 (redes sem fio, mobilidade e QoS)
    \item \textbf{LoRa Alliance} --- \emph{LoRaWAN 1.0.4 Specification}, 2020 (SF, tempo no ar, ADR)
    \item \textbf{IEEE Std 802.11ax-2021} e \textbf{IEEE Std 802.11be-2024} (OFDMA, MLO)
    \item \textbf{3GPP TS 38.300} (NR) e \textbf{3GPP TR 38.811} (estudo NTN)
    \item \textbf{Stallings}, \emph{Wireless Communications \& Networks} (LPWAN, 5G, satélites)
  \end{itemize}
  \begin{alertblock}{Dica de estudo}
    Refazer o cálculo de tempo no ar do exercício para $SF=7$/$BW=125$~kHz e comparar com $SF=12$: observe a redução drástica do tempo de transmissão.
  \end{alertblock}
\end{frame}

\end{document}